\subsection{ENVELOPING ALGEBRA OF THE OPPOSITE LIE ALGEBRA}

Let \(g\) be a Lie algebra over \(K\) , \(g^0\) the opposite Lie algebra and \(\sigma\) and \(\sigma_0\) the canonical mappings of \(g\) and \(g^0\) into their enveloping algebras \(U\) and \(V\) . Then \(\sigma\) is an \(\alpha\) -mapping of \(g^0\) into the associative algebra \(U^0\) opposite to the associative algebra \(U\) . Hence there exists one and only one homomorphism \(\phi\) of \(V\) into \(U^0\) mapping 1 to 1 and such that \(\sigma = \phi \circ \sigma_0\) .

PROPOSITION 4. The homomorphism \(\phi\) is an isomorphism of \(V\) onto \(U^0\) .

There exists a homomorphism \(\phi'\) of U into \(V^0\) mapping 1 to 1 and such that \(\sigma_0 = \phi' \circ \sigma\) . \(\phi'\) can be considered as a homomorphism of \(U^0\) into V. Then \(\sigma_0 = \phi' \circ \phi \circ \sigma_0\) and \(\sigma = \phi \circ \phi' \circ \sigma\) and hence \(\phi' \circ \phi\) and \(\phi \circ \phi'\) are the identity mappings of V and U. Hence the proposition.

V is identified with \(U^{0}\) under the isomorphism \(\phi\) . Then \(\sigma_{0}\) is identified with \(\sigma\) .

With this identification, the isomorphism \(\theta: x \mapsto -x\) of \(\mathfrak{g}\) onto \(\mathfrak{g}^0\) defines an isomorphism \(\tilde{\theta}\) of U onto V = U \(^0\) . This isomorphism can be considered as an antiautomorphism of U. It is called the principal antiautomorphism of U. If \(x_{1}, \ldots, x_{n}\) are in \(\mathfrak{g}\) , then:

\[
\begin{array}{r l} \tilde {\theta} (\sigma (x _ {1}) \dots \sigma (x _ {n})) & = \tilde {\theta} (\sigma (x _ {n})) \dots \tilde {\theta} (\sigma (x _ {1})) = (- \sigma (x _ {n})) \dots (- \sigma (x _ {1})) \\ & = (- 1) ^ {n} \sigma (x _ {n}) \dots \sigma (x _ {1}). \end{array}\tag{2}
\]

